% One offset suffices when the target coefficients absorb a possible borrow.
\section{Allowing a borrow in the coefficient tests}\label{sec:borrow-mask}

The low part of the coefficient-testing offset can be removed. The
resulting digits may receive a borrow from the preceding position, but
the residuals can be encoded so that both possible digits give the same
answer. Starting from Theorem~\ref{thm:best104}, replace
\[
 S_3=B\lambda(1+q)-(Z\lambda-2e)\mathcal C^2
\]
by
\begin{equation}\label{eq:borrow-third-block}
 S_3=B\lambda q-(Z\lambda-2e)\mathcal C^2,
 \qquad \sigma=S_3>0.
\end{equation}
Here $\mathcal C=x+g$ is the digit code, distinguished from the main
Pell denominator $c$. The supplied positive witness $\sigma$ remains
part of the system. The change removes the instruction computing
$q+1$. It requires a new fixed coefficient index, constructed below.

\subsection{Coefficient targets and the admissible index}

Keep the modular support of Section~\ref{sec:modular-support}, including
the input at weight zero, $\delta$ at $v_0=1$, the original 58
coordinates at $v_1,\ldots,v_{58}$ and the guard coordinate $u$ at
$v_{59}$. In particular,
\[
\begin{gathered}
 v_i=1+3\bigl(122i+(i^2\bmod61)\bigr),\qquad 0\le i\le59,\\
 v_*=21607,\qquad t_j=158467571+86429j\quad(0\le j<1832),\\
 t_*=316719070,\qquad K=t_*+1=316719071,\qquad
 L=2^{32}>3K+2.
\end{gathered}
\]
The 1830 original homogeneous quadratic residuals come from the
integer row basis and padding of Section~\ref{sec:quadratic-masks}.
Adjoin the guard residual $u\delta-x^2$ as the 1831st ordinary row.
For each ordinary row replace its coefficient target by
\begin{equation}\label{eq:borrow-targets}
 G_j=4F_j^h+\delta^2\qquad(0\le j<1831),
 \qquad G_{1831}=\delta^2.
\end{equation}
Thus $G_{1830}=4(u\delta-x^2)+\delta^2$, and the last row remains
the special square test.

For a homogeneous quadratic monomial $I$, let $c_I\in\{1,2\}$ be its
coefficient in $\mathcal C^2$: it is 1 for a square and 2 for a
product of distinct coordinates. If $a_{j,I}$ is its coefficient in
$G_j$, define
\[
 \mathcal D(B)=\sum_{j,I}\frac{a_{j,I}}{c_I}
                   B^{t_j-w(I)}.
\]
Every displayed coefficient is integral. The coefficients from
$4F_j^h$ are divisible by both possible denominators, and the added
$\delta^2$ is a square with denominator 1. In the special row the
coefficient 1 is placed at $t_{1831}-2$.

The support separation proved in Section~\ref{sec:modular-support}
gives the exact identity
\begin{equation}\label{eq:borrow-coefficient-identity}
 [B^{t_j}]\,\mathcal D(B)\mathcal C(B)^2=G_j.
\end{equation}
Indeed, distinct monomials have distinct weights, different row
intervals are disjoint, and a product involving a dummy row coordinate
lies above every row target. The least position of $\mathcal D$ is
above $v_*$, so all low coordinate tests remain automatic.

Choose a power of two $z\ge2$ strictly above the absolute values of
the new coefficients of $\mathcal D$, and put $Z=2z$. Define
\begin{align*}
 \ell_0(B)&=\sum_{i=0}^{59}B^{v_i}+\sum_{j=0}^{1831}B^{t_j},\\
 e_0(B)&=z\sum_{j=0}^{K-1}B^j+\mathcal D(B),\\
 V&=\ell_0(Z)+e_0(Z)Z^L.
\end{align*}
All digits of $e_0$ below $K$ lie in $[1,Z-1]$; higher digits vanish.
Choose a power of two
\begin{equation}\label{eq:borrow-index-bound}
 H>\max\{2Z^{2L+1},\ 4^{t_*+3}Z\,1900^2,\ 3L,\ 16\}.
\end{equation}
The positions, coordinate count and size formula for $H$ are unchanged.
The indices $Z,V,H$ are rebuilt for the new coefficient polynomial.
They depend effectively on the represented set and are independent of
the queried positive input $x$.

\subsection{Bounds before the coefficients are decoded}

Retain the coding equations and positive conditions
\begin{align*}
 B&=Hb^2,& b&=x+\beta,& \lambda(B-1)&=q^2-1,& \theta&=B-Z,\\
 S_2&=\ell+eq=V+t\theta,& S_2+\alpha&=q^2,&
 \Omega&=Z\lambda-2e>0,& \mathcal C&=x+g.
\end{align*}
The packing remains
\begin{align*}
 n&=q^8,& S&=g+q^2(S_2+q^2S_3),\\
 T+1&=q^2(1+\theta\lambda)-(b-1)\ell+(B-4)\ell q^4,\\
 r&=S(n^2-n)+(T+1)(n^2-1).
\end{align*}
The Pell interval scale $Y_P=sn^2$ is distinct from the packed code
$S_2$.

All preliminary upper bounds from Sections~\ref{sec:final-two-savings}
and~\ref{sec:unit-scale} still apply, since the new positive offset
$B\lambda q$ is smaller than $B\lambda(1+q)$. More explicitly,
$S_2<q^2$ gives $e<q$ and $\ell<q^2$, while the geometric equation
gives $B\le q^2$. Positivity of $\Omega$ and $\sigma$ gives
\[
 \mathcal C^2<B\lambda q<3q^3<q^4,
 \qquad 0<S_3<q^4.
\]
Thus $g<q^2$ and the same normalized masks give
\[
 0<S<n,\qquad0\le T<n,\qquad
 n\le r\le2n^3-2n^2<2n^3.
\]
The Pell proof of Theorem~\ref{thm:best104} uses these bounds and the
positive interval, independently of the coefficient values in the
encoding. It therefore recovers $b$ as a power of two, $q=B^L$, and
$n^2\mid\binom{2r}{r}$.

After $q=B^L$ is known, $\lambda>q$ and $e<q$ imply
$\Omega>Z\lambda/2$. Equation~\eqref{eq:borrow-third-block} then gives
\begin{equation}\label{eq:borrow-code-bound}
 \mathcal C^2<\frac{2Bq}{Z}\le\frac{Bq}{2}<q^2.
\end{equation}
Hence $g<\mathcal C<q$ before either the input code or the quotient
of the packed polynomial is decoded.

The first two masks and the packed congruence have not changed.
The borrow-removal and base-$Z$ uniqueness arguments of
Sections~\ref{sec:reversed-packing} and~\ref{sec:input-unit} require
only the bounds just proved, the support of $\ell_0$, the digit bounds
for $e_0$, and~\eqref{eq:borrow-index-bound}. They therefore give
\begin{equation}\label{eq:borrow-alias}
 S_2=\ell_0(B)+e_0(B)q,\qquad
 \ell=\ell_0(B)+mq,\qquad e=e_0(B)-m,\qquad 0\le m<e_0(B).
\end{equation}
The first mask, using $g<q$, decodes the input digit exactly as $x$.
All other coordinates are nonnegative integers less than $b$; dummy
coordinates can occur only at the row positions. Consequently
\begin{equation}\label{eq:borrow-polynomial-bounds}
 \mathcal C<B^K,\qquad
 A_C:=\mathcal C(1)^2<(1900b)^2,\qquad
 0<e<B^K,\qquad 2e\mathcal C^2<q.
\end{equation}
The last estimate uses $L>3K+2$ and is valid before $m$ is excluded.

\subsection{The high digits still exclude the quotient ambiguity}

\begin{lemma}\label{lem:borrow-high-window}
Under the preceding hypotheses, the third no-carry condition forces
$m=0$ in~\eqref{eq:borrow-alias}.
\end{lemma}
\begin{proof}
Put $J_C=ZA_C<B/4$ and
\[
 W_C=Z\lambda\mathcal C^2=qW_1+W_0,
 \qquad0\le W_0<q.
\]
The coefficients of the polynomial $Z\lambda\mathcal C^2$ are
nonnegative and at most $J_C$. They are already its base-$B$ digits,
so no internal carries occur in this product. Since
$\deg\mathcal C^2\le2K-2$, these coefficients equal $J_C$ throughout
the interval $[2K-2,2L-1]$. In particular they equal $J_C$ at every
position $L,\ldots,L+K$, because $L>3K+2$.

By~\eqref{eq:borrow-polynomial-bounds}, exact division by $q$ gives
\[
 \left\lfloor\frac{S_3}{q}\right\rfloor
  =B\lambda-W_1+\varepsilon,
 \qquad
 \varepsilon=\left\lfloor
   \frac{2e\mathcal C^2-W_0}{q}\right\rfloor\in\{-1,0\}.
\]
The unit digit of $B\lambda$ is zero, and its digits at positions
$1,\ldots,2L$ are one. The first $K+1$ digits of $W_1$ are all
$J_C$. At the unit position the subtraction therefore gives digit
$B-J_C-1$ or $B-J_C$ and an outgoing borrow of one. At each following
position through $K$, the incoming borrow changes $1-J_C$ into
$-J_C$, giving digit $B-J_C$ and again an outgoing borrow of one.
Thus the digits of $S_3$ at positions $L,\ldots,L+K$ are all at least
$B-J_C-1$.

The product $(B-4)\ell_0(B)$ has digits $B-4$ at its support
positions and no internal carries, and it is less than $q$.
Accordingly the part of the third mask at and above position $L$ is
\[
 X=(B-4)m<B^{K+1}.
\]
If a digit $d$ has no binary overlap with a digit at least
$B-J_C-1$, then
$d\le B-1-(B-J_C-1)=J_C$.
Hence every digit of $X$ is at most $J_C$.

Let $F$ be the digit polynomial of $X$, so $F(B)=X$, its coefficients
are nonnegative, and $\deg F\le K$. Since $B-4\mid F(B)$, reduction
modulo $B-4$ gives $B-4\mid F(4)$. On the other hand,
\[
 0\le F(4)\le J_C\frac{4^{K+1}-1}{3}
       <\frac{B}{12}<B-4.
\]
Here~\eqref{eq:borrow-index-bound} and
$J_C<Z(1900b)^2$ give
$4^{K+2}J_C<B$, because $K=t_*+1$.
It follows that $F(4)=0$. Nonnegativity of its coefficients then
gives $F=0$, so $X=0$ and $m=0$.
\end{proof}

\subsection{Low carries determine the intended equations}

We now have $\ell=\ell_0(B)$ and $e=e_0(B)$.
At positions below $K$, the offset $qB\lambda$ makes no contribution,
and the negative tail of $e_0-z\lambda$ begins too high to contribute.
The raw coefficients there are precisely those of
$2\mathcal D(B)\mathcal C(B)^2$. Each has absolute value less than
$2zA_C=J_C<B/4$.

Normalize these signed coefficients starting at the unit position.
The initial carry is zero. If an incoming carry is in $\{-1,0\}$,
adding a raw coefficient of absolute value less than $B/4$ gives a
number strictly between $-B$ and $B$. Its Euclidean quotient by $B$
is therefore again -1 or 0. Induction proves that these are the only
incoming carries at every target position.

At a target, let $v$ be the raw coefficient and let
$\varepsilon\in\{-1,0\}$ be the incoming carry. The mask digit
$B-4$ permits exactly the digits $0,1,2,3$. Since $|v|<B/4$, a
negative $v+\varepsilon$ would normalize to a digit greater than 3,
and a nonnegative one cannot wrap around $B$. Thus the target test is
exactly
\begin{equation}\label{eq:borrow-local-test}
 0\le v+\varepsilon\le3.
\end{equation}
At the special row, $v=2\delta^2$. The test implies
$2\delta^2\le4$, so the nonnegative integer $\delta$ is 0 or 1.
At the guard, the raw value is
\[
 v=8(u\delta-x^2)+2\delta^2.
\]
If $\delta=0$, this is $-8x^2<0$, and a carry of -1 or 0 cannot
make it satisfy~\eqref{eq:borrow-local-test}. Hence $\delta=1$.

Every ordinary row now has $v=8F_j^h+2$. Depending on the carry,
\eqref{eq:borrow-local-test} says
\[
 0\le8F_j^h+1\le3
 \quad\hbox{or}\quad
 0\le8F_j^h+2\le3.
\]
Both imply $F_j^h=0$ by integrality. Since $\delta=1$, these are
exactly the original equations. At each low variable position the
raw coefficient and incoming carry are zero, since the least position
of $\mathcal D$ lies above all such positions. Those tests impose no
additional condition. This proves soundness of the new encoding.

\subsection{Positive witnesses and the resulting certificate}

Conversely, start with a solution of the represented system. Set
$\delta=1$, $u=x^2$, and every dummy coordinate to zero. Choose a
sufficiently large power of two $b$, and take the canonical
$B,q,\lambda,\ell_0,e_0,\mathcal C$ defined above. The positive digit
at weight one makes $g>0$. All previous fixed-index, packed-congruence,
bound and $\Omega$ witnesses remain positive, because the support and
digit bounds have the same form.

The changed witness $\sigma=S_3$ is positive as well. Indeed,
$\mathcal C<B^K$, $Z<B$ and $L>2K$ give
\[
 Z\mathcal C^2<B^{2K+1}<qB,
 \qquad
 S_3=\lambda(qB-Z\mathcal C^2)+2e\mathcal C^2>0.
\]
At every ordinary row and the special row the raw coefficient is
exactly 2. Its actual digit is consequently 1 or 2, both allowed by
the mask. The low coordinate tests have digit zero, and $\ell_0$ has
no other support. The first two masks work as before. Thus the
packed integers have no binary overlap, yielding
$n^2\mid\binom{2r}{r}$ for their newly calculated $r$.

All dependent positive Pell witnesses are chosen by the constructions
of Section~\ref{sec:unit-scale}, using this $r$. Those constructions
require only the packed ranges, the binomial divisibility, and the
fixed exponent bounds, which have just been established. This proves
necessity without identifying the old and new values of $r$.

\begin{theorem}\label{thm:best103}
For each r.e.\ set, the modified coefficient targets
\eqref{eq:borrow-targets} and admissible index above give a
membership-equivalent system in 34 positive unknowns and 22 equations
with a straight-line certificate of \textbf{103 arithmetic instructions}:
57 multiplications and 46 additions. Fixed numerals, supplied parameters
and equality tests are free.
\end{theorem}
\begin{proof}
The preceding arguments prove soundness and positive necessity.
The product $\mathit{Lam}=B\lambda$ is already required by the
geometric equation. Replace the instruction
$P_2=\mathit{Lam}(q+1)$ by $P_2=\mathit{Lam}\,q$ and delete the
single addition computing $q+1$. All other instruction counts remain
unchanged. The coefficient adjustment changes only the supplied fixed
index; it introduces no system unknown or arithmetic operation.
The exact primitive and source-residual identities, including the two
inherited triangular corrections, are checked by
\file{round20\_1980\_borrow\_mask\_certificate.py}.
The full high-window argument is Lemma~\ref{lem:borrow-high-window};
it needs no stronger bound on $H$ than the preceding construction.
\end{proof}
